% GENERATED FILE. Edit canonical lessons, then run npm run generate:books.
% canonical-source-hash: fnv1a64:1e8caf711b877728
% canonical-lessons: GU-C238-biladi, GU-C238-garoli, GU-C238-karoliyo, GU-C238-kidi, GU-C238-latakvu

\chapter{Creatures in the house}
\label{ch:gu-creatures-in-the-house}
\hlchaptermodality{238}

\begin{chapteropening}
\textbf{By the end of this chapter:} \emph{I can say what the small creatures in the house did.}

Complete the last lesson of chapter 238: I can say what the small creatures in the house did.
\end{chapteropening}

\section[\texorpdfstring{\gu{બિલાડી} (bilāḍī)}{બિલાડી (bilāḍī)}]{\gu{બિલાડી} (bilāḍī) — cat}

\label{lesson:GU-C238-biladi}



\begin{quote}
Before the new one: say the Gujarati for to flutter, then the Gujarati for to fly up.
\end{quote}

\subsection*{You'll want to know: \gu{બિલાડી}}
\textbf{\gu{બિલાડી}} — \emph{bilāḍī} — \textquotedblleft{}cat\textquotedblright{}.

\begin{quote}
\textbf{\gu{બિલાડી} \gu{રાત્રે} \gu{ઘરમાં} \gu{આવી}.} — \emph{bilāḍī rātre gharmā̃ āvī} — \textquotedblleft{}The cat came into the house at night.\textquotedblright{}
\end{quote}

\begin{cousinweb}
\textbf{\gu{બિલાડી}} (\emph{bilāḍī}) is a tadbhava, a word inherited from Sanskrit and reshaped by sound change. It comes from \emph{biḍāla}, \textquotedblleft{}cat\textquotedblright{} (CDIAL 9237).
\end{cousinweb}

\begin{grammarlens}[title={What you've built}]
One more word for everyday life at A2.
\end{grammarlens}

\subsection*{Your turn}
\begin{itemize}
\raggedright
  \item \emph{Say it:} \emph{bilāḍī}
  \item \emph{Say it:} \emph{bilāḍī}, once more
  \item \emph{Say it:} \emph{bilāḍī rātre gharmā̃ āvī} — The cat came into the house at night.
  \item \emph{Recall:} say the Gujarati for to flutter, then the Gujarati for to fly up, then say \emph{bilāḍī} again
\end{itemize}

\begin{tcolorbox}[breakable,colback=teal!4,colframe=teal!35!black,title={Before you move on}]
What does \gu{બિલાડી} mean? (\textquotedblleft{}Cat\textquotedblright{}.) Say it once more.
\end{tcolorbox}
\section[\texorpdfstring{\gu{ગરોળી} (garoḷī)}{ગરોળી (garoḷī)}]{\gu{ગરોળી} (garoḷī) — house lizard, gecko}

\label{lesson:GU-C238-garoli}



\begin{quote}
Before the new one: say the Gujarati for to fly up, then the Gujarati for cat.
\end{quote}

\subsection*{You'll want to know: \gu{ગરોળી}}
\textbf{\gu{ગરોળી}} — \emph{garoḷī} — \textquotedblleft{}house lizard, gecko\textquotedblright{}.

\begin{quote}
\textbf{\gu{ગરોળી} \gu{બારી} \gu{પાસે} \gu{હતી}.} — \emph{garoḷī bārī pāse hatī} — \textquotedblleft{}The lizard was near the window.\textquotedblright{}
\end{quote}

\begin{cousinweb}
\textbf{\gu{ગરોળી}} (\emph{garoḷī}) is inherited. CDIAL lists the Gujarati \emph{gharoḷī}, \textquotedblleft{}house lizard\textquotedblright{}, under \emph{*gharagolī}, a reconstructed form built on \emph{ghara}, \textquotedblleft{}house\textquotedblright{} (CDIAL 4431).
\end{cousinweb}

\begin{grammarlens}[title={What you've built}]
One more word for everyday life at A2.
\end{grammarlens}

\subsection*{Your turn}
\begin{itemize}
\raggedright
  \item \emph{Say it:} \emph{garoḷī}
  \item \emph{Say it:} \emph{garoḷī}, once more
  \item \emph{Say it:} \emph{garoḷī bārī pāse hatī} — The lizard was near the window.
  \item \emph{Recall:} say the Gujarati for to fly up, then the Gujarati for cat, then say \emph{garoḷī} again
\end{itemize}

\begin{tcolorbox}[breakable,colback=teal!4,colframe=teal!35!black,title={Before you move on}]
What does \gu{ગરોળી} mean? (\textquotedblleft{}House lizard, gecko\textquotedblright{}.) Say it once more.
\end{tcolorbox}
\section[\texorpdfstring{\gu{કરોળિયો} (karoḷiyo)}{કરોળિયો (karoḷiyo)}]{\gu{કરોળિયો} (karoḷiyo) — spider}

\label{lesson:GU-C238-karoliyo}



\begin{quote}
Before the new one: say the Gujarati for cat, then the Gujarati for house lizard.
\end{quote}

\subsection*{You'll want to know: \gu{કરોળિયો}}
\textbf{\gu{કરોળિયો}} — \emph{karoḷiyo} — \textquotedblleft{}spider\textquotedblright{}.

\begin{quote}
\textbf{\gu{છત} \gu{પાસે} \gu{એક} \gu{કરોળિયો} \gu{હતો}.} — \emph{chat pāse ek karoḷiyo hato} — \textquotedblleft{}There was a spider near the ceiling.\textquotedblright{}
\end{quote}

\begin{cousinweb}
\textbf{\gu{કરોળિયો}} (\emph{karoḷiyo}) is a tadbhava. CDIAL lists it, glossed \textquotedblleft{}spider\textquotedblright{}, under Sanskrit \emph{kolika}, \textquotedblleft{}weaver\textquotedblright{}, so the spider is named for its weaving (CDIAL 3535).
\end{cousinweb}

\begin{grammarlens}[title={What you've built}]
One more word for everyday life at A2.
\end{grammarlens}

\subsection*{Your turn}
\begin{itemize}
\raggedright
  \item \emph{Say it:} \emph{karoḷiyo}
  \item \emph{Say it:} \emph{karoḷiyo}, once more
  \item \emph{Say it:} \emph{chat pāse ek karoḷiyo hato} — There was a spider near the ceiling.
  \item \emph{Recall:} say the Gujarati for cat, then the Gujarati for house lizard, then say \emph{karoḷiyo} again
\end{itemize}

\begin{tcolorbox}[breakable,colback=teal!4,colframe=teal!35!black,title={Before you move on}]
What does \gu{કરોળિયો} mean? (\textquotedblleft{}Spider\textquotedblright{}.) Say it once more.
\end{tcolorbox}
\section[\texorpdfstring{\gu{કીડી} (kīḍī)}{કીડી (kīḍī)}]{\gu{કીડી} (kīḍī) — ant}

\label{lesson:GU-C238-kidi}



\begin{quote}
Before the new one: say the Gujarati for house lizard, then the Gujarati for spider.
\end{quote}

\subsection*{You'll want to know: \gu{કીડી}}
\textbf{\gu{કીડી}} — \emph{kīḍī} — \textquotedblleft{}ant\textquotedblright{}.

\begin{quote}
\textbf{\gu{કીડીઓ} \gu{ગોળ} \gu{પાસે} \gu{આવી}.} — \emph{kīḍīo goḷ pāse āvī} — \textquotedblleft{}The ants came to the jaggery.\textquotedblright{}
\end{quote}

\begin{cousinweb}
\textbf{\gu{કીડી}} (\emph{kīḍī}) is a tadbhava. It comes from Sanskrit \emph{kīṭa}, \textquotedblleft{}insect, worm\textquotedblright{}, and CDIAL glosses the Gujarati word \textquotedblleft{}ant\textquotedblright{} (CDIAL 3193).
\end{cousinweb}

\begin{grammarlens}[title={What you've built}]
One more word for everyday life at A2.
\end{grammarlens}

\subsection*{Your turn}
\begin{itemize}
\raggedright
  \item \emph{Say it:} \emph{kīḍī}
  \item \emph{Say it:} \emph{kīḍī}, once more
  \item \emph{Say it:} \emph{kīḍīo goḷ pāse āvī} — The ants came to the jaggery.
  \item \emph{Recall:} say the Gujarati for house lizard, then the Gujarati for spider, then say \emph{kīḍī} again
\end{itemize}

\begin{tcolorbox}[breakable,colback=teal!4,colframe=teal!35!black,title={Before you move on}]
What does \gu{કીડી} mean? (\textquotedblleft{}Ant\textquotedblright{}.) Say it once more.
\end{tcolorbox}
\section[\texorpdfstring{\gu{લટકવું} (laṭakvũ)}{લટકવું (laṭakvũ)}]{\gu{લટકવું} (laṭakvũ) — to hang, to dangle}

\label{lesson:GU-C238-latakvu}



\begin{quote}
Before the new one: say the Gujarati for spider, then the Gujarati for ant.
\end{quote}

\subsection*{You'll want to know: \gu{લટકવું}}
\textbf{\gu{લટકવું}} — \emph{laṭakvũ} — \textquotedblleft{}to hang, to dangle\textquotedblright{}.

\begin{quote}
\textbf{\gu{કરોળિયો} \gu{હવામાં} \gu{લટક્યો}.} — \emph{karoḷiyo havāmā̃ laṭakyo} — \textquotedblleft{}The spider hung in the air.\textquotedblright{}
\end{quote}

\begin{cousinweb}
\textbf{\gu{લટકવું}} (\emph{laṭakvũ}) comes from \emph{*laṭṭa}, \textquotedblleft{}hanging, tangled\textquotedblright{}, a reconstructed form; CDIAL glosses the Gujarati verb \textquotedblleft{}to hang down\textquotedblright{} (CDIAL 10918).
\end{cousinweb}

\begin{grammarlens}[title={What you've built}]
One more word for everyday life at A2.
\end{grammarlens}

\subsection*{Your turn}
\begin{itemize}
\raggedright
  \item \emph{Say it:} \emph{laṭakvũ}
  \item \emph{Say it:} \emph{laṭakvũ}, once more
  \item \emph{Say it:} \emph{karoḷiyo havāmā̃ laṭakyo} — The spider hung in the air.
  \item \emph{Recall:} say the Gujarati for spider, then the Gujarati for ant, then say \emph{laṭakvũ} again
\end{itemize}

\begin{tcolorbox}[breakable,colback=teal!4,colframe=teal!35!black,title={Before you move on}]
What does \gu{લટકવું} mean? (\textquotedblleft{}To hang, to dangle\textquotedblright{}.) Say it once more.
\end{tcolorbox}
